\appendix
\titleformat{\section}[block]{\centering\bfseries\fontsize{20pt}{24pt}\selectfont}{ภาคผนวก \thesection}{0.6em}{}
\section{สรุปแบบจำลองและหลักฐานจากโค้ด}
\begin{figure}[H]
\centering
\begin{tikzpicture}[node distance=9mm, every node/.style={draw,rounded corners,align=center,font=\normalsize,text width=38mm,minimum height=12mm},>=Stealth]
\node (source) at (0,0) {ข้อมูลต้นทาง\\รายการและข้อมูลอ้างอิง};
\node (etl) at (5.2,0) {ETL\\สกัด แปลง ตรวจ และโหลด};
\node (dw) at (10.4,0) {คลังข้อมูลรุ่นใหม่\\มิติ Fact และ Bridge};
\node (api) at (10.4,-2.2) {บริการรายงาน\\ตรวจตัวตนและขอบเขต};
\node (adapter) at (5.2,-2.2) {ตัวแปลงข้อมูลรายงาน\\สรุปและจัดกลุ่ม};
\node (web) at (0,-2.2) {มุมมอง 6 บทบาท\\กราฟ ตาราง รายละเอียด};
\draw[->] (source)--(etl); \draw[->] (etl)--(dw); \draw[->] (dw)--(api);
\draw[->,dashed] (api)--(adapter); \draw[->] (adapter)--(web);
\end{tikzpicture}
\caption{สถาปัตยกรรมเป้าหมายของรายงาน เส้นประคือการเชื่อมที่ต้องยืนยันเพิ่มเติม เนื่องจาก provider ปัจจุบันใช้ข้อมูลจำลอง ภาพนี้วาดด้วยเวกเตอร์ในเอกสาร}
\end{figure}

\begin{figure}[H]
\centering
\begin{tikzpicture}[every node/.style={draw,rounded corners,align=center,font=\normalsize,text width=39mm,minimum height=13mm},>=Stealth]
\node (request) at (0,0) {Fact: ใบลา\\1 แถวต่อใบลา};
\node (student) at (-5,0) {มิตินิสิตและบุคลากร\\ผู้ลาและที่ปรึกษา};
\node (org) at (5,0) {มิติองค์กรและประเภทลา\\คณะ ภาค หลักสูตร};
\node (term) at (0,2.2) {มิติภาคการศึกษา\\ปี เทอม และลำดับ};
\node (rs) at (0,-2.2) {Fact: ใบลา--กลุ่มเรียน\\1 แถวต่อคู่คีย์};
\node (offering) at (-5,-2.2) {มิติรายวิชาและ Offering\\วิชา เทอม กลุ่มเรียน};
\node (bridge) at (5,-2.2) {ตารางเชื่อมผู้สอน\\รองรับหลายผู้สอน};
\draw[->] (student)--(request); \draw[->] (org)--(request); \draw[->] (term)--(request);
\draw[->] (request)--node[right,draw=none,text width=12mm,minimum height=0mm]{1:N}(rs);
\draw[->] (offering)--(rs); \draw[->] (rs)--(bridge);
\end{tikzpicture}
\caption{แบบจำลองเชิงแนวคิดเฉพาะการวิเคราะห์การลา ลูกศรแสดงการเชื่อมข้อมูล ไม่ใช่ ERD ครบทุกคีย์}
\end{figure}


\subsection{ตารางจับคู่คำถามกับโครงสร้างข้อมูล}
\renewcommand{\arraystretch}{1.3}
\noindent
\begin{minipage}{\linewidth}
  \captionof{table}{คำถามและขอบเขตข้อมูลตามบทบาทผู้ใช้}
  \begin{tabularx}{\textwidth}{p{25mm}XX}\toprule\rowcolor{tableheader}
    \textbf{บทบาท} & \textbf{คำถามหลัก} & \textbf{ข้อมูลและขอบเขต}\\\midrule
    นิสิต & ตนเองลาเท่าใด แต่ละเทอมและรายวิชาต่างกันอย่างไร & Request fact และ request--offering fact; เฉพาะตนเอง\\\midrule
    ที่ปรึกษา & ชั้นปีใดและนิสิตคนใดมีการลาสูง & Request fact, student, term; รายชื่อนิสิตที่ดูแลปัจจุบัน\\\midrule
    ผู้สอน & วิชาและกลุ่มเรียนใดได้รับผลกระทบ & Request--offering fact, offering, lecturer bridges; กลุ่มเรียนที่มีสิทธิ์\\\midrule
    หัวหน้าภาค & ภาพรวมนิสิตและวิชาของภาคเป็นอย่างไร & Student department และ subject owner department แยกกัน\\\midrule
    ผู้บริหาร & ภาค หลักสูตร และกลุ่มที่ปรึกษาต่างกันอย่างไร & Request fact กับมิติองค์กร; ภายในคณะที่มีสิทธิ์\\\midrule
    เจ้าหน้าที่ & กลุ่มที่ปรึกษาและชั้นปีใดควรติดตาม & Request fact, person, student; ภาควิชาของบัญชี\\\bottomrule
  \end{tabularx}
\end{minipage}

\subsection{เงื่อนไขอ่านตัวเลข}
จำนวนใบลาต้องตัดซ้ำด้วยคีย์ใบลา จำนวนผู้ลาตัดซ้ำด้วยคีย์นิสิต วันลารวมจาก fact ระดับใบลา ส่วนครั้งเรียนและชั่วโมงเรียนใช้ fact ระดับใบลา--offering ตามความครบถ้วนของข้อมูล เทอมยอดศูนย์ต้องอยู่ในฐานคำนวณค่าเฉลี่ยที่เลือก และการจัดอันดับทุกชุดต้องระบุปีและภาคการศึกษา

\subsection{ขอบเขตหลักฐาน}
เอกสารนี้อ้างอิงไฟล์ในเครื่อง ณ วันที่ตรวจ ไม่ได้ตรวจฐานข้อมูลที่ติดตั้งหรือรันชุดทดสอบแอปพลิเคชัน ข้อมูลจำลองอธิบายรูปแบบหน้ารายงานได้ แต่ใช้ยืนยันปริมาณการลาจริง ความเร็ว หรือความถูกต้องของยอดจริงไม่ได้ เส้นทางหลักฐานแบบละเอียดอยู่ใน docs/SOURCE-NOTES.md ของเอกสารฉบับนี้
